\section{Network Architecture and Correctness}\label{sec:network}

This section reconstructs the network from repository evidence and live cross-checks
rather than trusting any single diagram. The EX3400 core switch \emph{was} inspected
live during this review after the operator supplied password credentials (E-22, E-25);
that inspection produced one significant new finding and corrected two documented
facts. One constraint remains: OPNsense was deliberately not touched because no
read-only API credential was used (E-23). Items that rest only on repository
documentation, not on a live reading, are marked \emph{inferred}. See
\code{diagrams/logical-topology.mmd} and \code{diagrams/physical-topology.mmd}.

\begin{callout}
\textbf{Live switch inspection, 2026-07-23.} Hostname \code{KM-EX3400}, model
\code{ex3400-48p}, Junos 23.4R2-S7.4. Virtual-Chassis mode enabled with a single member
(FPC0 master). Both power supplies OK, CPU sensor 37 C, no active alarms. SNMP is
read-only and client-locked to \code{192.168.10.183/32} with \code{0.0.0.0/0 restrict}.
\code{bpdu-block-on-edge} is set, with edge on \code{ge-0/0/22}, \code{ge-0/0/36} and
\code{ge-0/0/41}.
\end{callout}

\subsection{VLAN and Addressing Design}

The estate runs seven VLANs on the EX3400 (ELS), gatewayed by OPNsense. The design is
a clean default-deny segmentation model with a dedicated out-of-band plane.

\begin{tabularx}{\linewidth}{@{}c p{2.3cm} p{3.4cm} X@{}}
\toprule
\rh{ID} & \rh{Name} & \rh{Subnet} & \rh{Purpose / population} \\
\midrule
\endhead
1  & mgmt/default & 192.168.10.0/24 & Proxmox hosts, corosync, services, switch mgmt (dense) \\
20 & trusted/OOB  & 192.168.20.0/24 & iDRAC x2 + IPMI (BMCs) + Ares jump leg \\
30 & servers      & 192.168.30.0/24 & NFS/PBS/GPU egress; dual-homed GPU and storage nodes \\
40 & IoT          & 192.168.40.0/24 & smart-home, sensors (sparse) \\
50 & VoIP         & 192.168.50.0/24 & CP-8841 phones, FreePBX (deferred) \\
60 & guest        & 192.168.60.0/24 & isolated guest WiFi \\
70 & lab          & 192.168.70.0/24 & CCNA/experimental; standalone \code{sandbox} node \\
\bottomrule
\end{tabularx}

The GPU and storage nodes (QuarkyLab, Jarvis, Randy) are dual-homed: management and
corosync stay on VLAN 1 while NFS, PBS, and internet egress ride VLAN 30. Corosync was
deliberately kept off VLAN 30 to hold the cluster ring on stable L2, which is sound
practice.

\subsection{EX3400 Port Map (live-verified, E-22 / E-25)}

\begin{tabularx}{\linewidth}{@{}p{2.4cm} p{4.2cm} p{2.2cm} X@{}}
\toprule
\rh{Port} & \rh{Attachment} & \rh{Mode} & \rh{VLANs} \\
\midrule
\endhead
ge-0/0/22 & Jarvis nic2 (onboard 1G) & Trunk & native 1 (+ tagged 30, now vestigial) \\
ge-0/0/24 & QuarkyLab nic2 & Trunk & native 1 + tagged 30 \\
ge-0/0/30, /32, /44 & iDRAC/IPMI OOB & Access & 20 only \\
ge-0/0/38 & APC AP7901 PDU & Access & 1 \\
ge-0/0/41 & Ares mgmt/jump & Dual & native 1 + tagged 20 \\
ge-0/0/45 & unused (no neighbor) & Down & none \\
ge-0/0/46 & UniFi USW-24 Port 24 & Trunk & native 1 + tagged 20/30/40/50/60/70 \\
xe-0/2/0 & Randy nic3 (10G ConnectX-3) & Trunk & 1 + 30 (storage/NFS) \\
xe-0/2/2 & Jarvis enp132s0 (10G ConnectX) & Access & 30 (dedicated storage/egress) \\
xe-0/2/3 & UniFi SFP2 & Up 10G (DESG) & native 1; link confirmed live (DR-14 resolved) \\
\bottomrule
\end{tabularx}

Two documented facts were corrected by the live inspection. First, \code{ge-0/0/45} does
not connect to an EX2300: the operator confirms \textbf{no such switch exists or ever
has}, and the port is simply down with no LLDP neighbor. The EX2300 was a documentation
phantom and has been removed from the inventory (DR-15). Second, \code{xe-0/2/3} is
\textbf{up at 10 Gbps} (media type copper, LLDP neighbor a Ubiquiti device) and forwards
as a designated port, so the previously recorded EEPROM speed-mismatch fault is stale and
the redundant UniFi path does exist (DR-14 resolved).

The ELS \code{native-vlan-id} question is now settled by live configuration rather than
by narrative: the switch carries \code{set interfaces ge-0/0/22 native-vlan-id 1} and the
same on \code{ge-0/0/24}, \code{ge-0/0/41}, \code{ge-0/0/46}, \code{xe-0/2/0} and
\code{xe-0/2/3}. It is set at the physical-interface
level, not under \code{unit 0 family ethernet-switching}. A stale note in two documents
still claims the EX3400 does not support it at all; that is false and would cause an
operator to build a broken trunk (drift DR-08, doc fix D-08).

\subsection{Gateway, DHCP, and DNS}

OPNsense (VM 100 on pve2) is the single LAN router, firewall, DHCP server, and
inter-VLAN gateway, live-confirmed running with high egress counters (E-06). DHCP moved
to the Kea backend after an empty-Kea outage on 2026-07-20; Kea holds 12 reservations
and serves all seven VLAN scopes. WAN is dual-path with single NAT: WAN1 is a Spectrum
public block terminated directly on \code{vtnet0} (avoiding double-NAT), WAN2 is a
FirstNet 5G hotspot on \code{vtnet2} as an automatic failover gateway group.

DNS is genuinely highly available: primary Pi-hole \code{.177} (pve1 CT103) and
secondary \code{.178} (pve5 CT108, a nebula-sync full mirror on a 15-minute timer).
OPNsense DHCP hands out both resolvers on every scope, so clients fail over
automatically. The resolver chain is client to Pi-hole to Unbound on \code{.1} to
internet. A residual concern is that both Pi-holes forward to a single Unbound instance
(enterprise finding R-03), which has historically been the root of intermittent DNS
behavior. Critical containers (\code{.177} Pi-hole, \code{.181} NPM, \code{.183}
Grafana, \code{.148} Homepage) were made static-in-container so a DHCP-server outage
cannot drop routing, a good hardening move.

\subsection{Firewall Policy}

The inter-VLAN posture is default-deny with explicit allows, and several clamps have
been applied and verified:

\begin{itemize}
\item \textbf{BMC isolation and egress clamp (verified 2026-07-17):} all three BMCs sit
on VLAN 20; only Ares reaches them; a floating quick-block prevents VLAN 20 to internet,
so compromised BMC firmware cannot reach VLAN 1, VLAN 30, or the internet.
\item \textbf{VLAN 30 to VLAN 1 block:} one surgical rule blocks servers to
192.168.10.0/24, with a narrow allow for Prometheus scrape on \code{:9100}.
\item \textbf{DNS pinholes:} isolated VLANs (IoT/VoIP/Guest/Lab) keep \code{:53} to the
Pi-hole alias while remaining isolated.
\item \textbf{Deny logging:} all eight VLAN deny rules log to Wazuh/Loki for
lateral-movement visibility.
\end{itemize}

\begin{callout}
\textbf{Standing decision respected - VLAN 1 egress lockdown is NOT enforced.} The
management-plane to internet limit is documented as deliberately deferred, and this
review does not recommend enforcing it without first satisfying every prerequisite:
UniFi discovery and required multicast/broadcast behavior, legitimate network-device
outbound connectivity, RKE2 and container-registry/CDN dependencies, Mist or other
vendor callbacks, administrative access paths, DHCP reservations, an emergency access
path, current firewall API permissions, and a tested rollback. Enforcing it prematurely
is a credible way to cause a self-inflicted outage.
\end{callout}

\subsection{Remote Access and Reverse Proxy}

Headscale (self-hosted, LXC 105 on pve5, v0.29.1) is the internal admin mesh only, with
nine nodes; students and researchers do not use it. Student and researcher ingress is a
Cloudflare Tunnel (\code{quarkylab.kylemason.org}) with no inbound ports and the server
IP hidden. Public web services are fronted by Nginx Proxy Manager (LXC 101, admin
\code{:81} firewalled to Ares only, F-05):

\begin{tabularx}{\linewidth}{@{}p{4.3cm} p{4.2cm} X@{}}
\toprule
\rh{Hostname} & \rh{Backend} & \rh{TLS / auth} \\
\midrule
\endhead
vault.kylemason.org & Vaultwarden \code{.30.182} & CF DNS-01 / login \\
grafana.kylemason.org & Grafana \code{.183} (pve4) & CF DNS-01 / Grafana \\
homepage.kylemason.org & Homepage \code{.148} & CF DNS-01 / basic auth \\
wazuh.kylemason.org & Wazuh \code{.184} & CF DNS-01 / Wazuh \\
chat.netframe.local & Open WebUI \code{.30.185} & internal / admin acct \\
llm.netframe.local & Jarvis llm\_router \code{.31:8000} & HTTP, Pi-hole-scoped DNS \\
\bottomrule
\end{tabularx}

\subsection{Spanning Tree: Root Was Elected by Accident (Now Remediated)}

The most substantive finding from the live inspection concerns spanning tree. The
documentation asserted that the EX3400 was the RSTP root at bridge-priority 4096. That
was \textbf{false}. The committed configuration contained \textbf{no
\code{bridge-priority} statement at all}, leaving the core switch at the default 32768.
The UniFi USW-24-250W (\code{b4:fb:e4:1b:ca:3c}) had therefore won the root election on
the lower MAC address, and the core reached it through \code{ge-0/0/46} at a root cost
of 20000 over a 1 Gb link.

This was not an outage, but it is a genuine L2 hardening gap: root placement was decided
by MAC accident rather than by design, so a reboot or topology change could relocate the
root unpredictably, and traffic paths were computed toward an access switch rather than
the central, dual-PSU core.

\begin{callout}
\textbf{Remediated during this review} (operator applied, reviewer verified, E-28).
\code{set protocols rstp bridge-priority 4096} was committed. The switch is now genuinely
root: Root ID and Bridge ID are both \code{4096.c0:42:d0:f9:ec:71}, and \code{ge-0/0/46}
changed role from ROOT to DESG. Worth recording: the first attempt auto-reverted because
a \code{commit confirmed} was never confirmed. The safety net behaved exactly as intended
and the change only persisted once deliberately committed. Tracked as QL-R23 / QL-REC21.
\end{callout}

\subsection{Findings}

\begin{itemize}
\item \sevhigh OPNsense is a single point of failure for all routing, DHCP, and
inter-VLAN traffic (QL-R04). A CARP HA plan exists (QL-REC10).
\item \sevmed No automated network-configuration capture (F-N1, R-08, QL-R16). The
switch config was captured once by hand during this review; key-based SSH is still not
configured and nightly capture is still absent. Deploy Oxidized (QL-REC09).
\item \sevmed Single Unbound recursion point on \code{.1} (enterprise R-03).
\item \sevlow A residual Amazon Echo/Fire tablet remains on flat VLAN 1 pending the IoT
cutover.
\item \sevinfo IPv6 is not configured on the VLANs and is default-denied; no explicit
jumbo/MTU tuning is documented on the 10G fabric.
\item \sevinfo \textbf{Resolved this review:} STP root pinned to the core (QL-R23);
\code{xe-0/2/3} confirmed up at 10G (DR-14); EX2300 removed as a phantom asset (DR-15);
\code{native-vlan-id} placement verified in live config (D-08).
\end{itemize}
